Of the \data{audit/sample_count_int} transcripts with the highest LLM-validated score, the human auditor rated \data{audit/true_positive_count_int} as containing collusive content and \data{audit/false_positive_count_int} as lacking collusive content for a true positive rate of \data{audit/true_positive_rate_pct_percentage1}. Table~\ref{tab:X2} reports the true positive rate after partitioning the sample into three roughly equal-sized subsamples (while respecting score breaks). The true positive rate is decreasing in the score with a true positive rate of \data{audit/top_bin/true_positive_rate_pct_percentage1} for the \data{audit/top_bin/count_int} transcripts in the highest bin and \data{audit/bottom_bin/true_positive_rate_pct_percentage1} for the \data{audit/bottom_bin/count_int} transcripts in the lowest bin.

With details provided later, 163 of the 192 false positive transcripts contained content that, while relevant, did not cross the line for facilitating coordinated conduct. The other 29 false positive transcripts lacked any relevant content (``no content'') in which case it may be appropriate to classify them as hallucinations. With that definition, the hallucination rate is 9.6\% (29 out of 301).

\input{../data/outputs/tables/human_audit_performance_bins.tex}

As detailed in the Supplementary Information, there are 71 distinct companies in these 109 true positive transcripts and the four companies with the largest number of transcripts are Gol Linhas Aéreas Inteligentes (7 transcripts) which is a domestic airline in Brazil, Micron Technology (6) which is a manufacturer of computer memory and data storage, Norske Skogindustrier (4) which is a pulp and paper manufacturer, and Coca-Cola Bottlers Japan (4). We will later take a closer look at Gol Linhas Aéreas Inteligentes.

22 of the 109 transcripts are from previously documented episodes of companies making public announcements with collusive content, including 15 associated with the airlines case. 10 of those 15 airlines transcripts were from companies in the \textcite{aryal2022coordinated} study: US Airways (3 transcripts), American Airlines (2), Delta Airlines (2), United Airlines (2), and AirTran (1); and then there were five transcripts from Air France-KLM (3) and Frontier Airlines (2). The other previously documented episodes are: auto rental (Avis), broiler chicken (Pilgrim’s Pride), generic drugs (Teva Pharmaceuticals), and pork (Smithfield Foods, WH Group).

In some of these cases – generic drugs for sure and possibly broiler chicken and pork – collusion also involved the more standard private communications. The LLM flagging them is a reminder that screening public communications for collusive content could lead to investigations that uncover the most compelling evidence of illegal collusion: direct, express, and private communications between competitors.
